\documentclass[10pt, letterpaper]{article}

% Packages:
\usepackage[
    ignoreheadfoot, % set margins without considering header and footer
    top=1.3 cm, % seperation between body and page edge from the top
    bottom=1.3 cm, % seperation between body and page edge from the bottom
    left=1.7 cm, % seperation between body and page edge from the left
    right=1.7 cm, % seperation between body and page edge from the right
    footskip=0.8 cm, % seperation between body and footer
    % showframe % for debugging
]{geometry} % for adjusting page geometry
\usepackage{titlesec} % for customizing section titles
\usepackage{tabularx} % for making tables with fixed width columns
\usepackage{array} % tabularx requires this
\usepackage{xcolor}
\definecolor{primaryColor}{RGB}{0, 0, 0} % define primary color
\usepackage{enumitem} % for customizing lists
\usepackage{fontawesome5} % for using icons
\usepackage[
    pdftitle={Dimitry's CV},
    pdfauthor={Dimitry Grebenyuk},
    colorlinks=true,
    urlcolor=primaryColor
]{hyperref} % for links, metadata and bookmarks
\usepackage[pscoord]{eso-pic} % for floating text on the page
\usepackage{calc} % for calculating lengths
\usepackage{changepage} % for one column entries (adjustwidth environment)
\usepackage{paracol} % for two and three column entries
\usepackage{ifthen} % for conditional statements
\usepackage{needspace} % for avoiding page brake right after the section title
\usepackage{iftex} % check if engine is pdflatex, xetex or luatex

% Ensure that generate pdf is machine readable/ATS parsable:
\ifPDFTeX
    \input{glyphtounicode}
    \pdfgentounicode=1
    \usepackage[T1]{fontenc}
    \usepackage[utf8]{inputenc}
    \usepackage{lmodern}
\fi

%\usepackage{charter}
\usepackage{mathpazo}

% Some settings:
\raggedright
\AtBeginEnvironment{adjustwidth}{\partopsep0pt} % remove space before adjustwidth environment
\pagestyle{empty} % no header or footer
\setcounter{secnumdepth}{0} % no section numbering
\setlength{\parindent}{0pt} % no indentation
\setlength{\topskip}{0pt} % no top skip
\setlength{\columnsep}{0.15cm} % set column seperation
\pagenumbering{gobble} % no page numbering

\titleformat{\section}{\needspace{4\baselineskip}\bfseries\large}{}{0pt}{}[\vspace{1pt}\titlerule]

\titlespacing{\section}{
    % left space:
    -1pt
}{
    % top space:
    0.2 cm
}{
    % bottom space:
    0.1 cm
} % section title spacing

\renewcommand\labelitemi{$\vcenter{\hbox{\small$\bullet$}}$} % custom bullet points
\newenvironment{highlights}{
    \begin{itemize}[
        topsep=0.05 cm,
        parsep=0.05 cm,
        partopsep=0pt,
        itemsep=0pt,
        leftmargin=0 cm + 10pt
    ]
}{
    \end{itemize}
} % new environment for highlights


\newenvironment{highlightsforbulletentries}{
    \begin{itemize}[
        topsep=0.05 cm,
        parsep=0.05 cm,
        partopsep=0pt,
        itemsep=0pt,
        leftmargin=10pt
    ]
}{
    \end{itemize}
} % new environment for highlights for bullet entries

\newenvironment{onecolentry}{
    \begin{adjustwidth}{
        0 cm + 0.00001 cm
    }{
        0 cm + 0.00001 cm
    }
}{
    \end{adjustwidth}
} % new environment for one column entries

\newenvironment{twocolentry}[2][]{
    \onecolentry
    \def\secondColumn{#2}
    \setcolumnwidth{\fill, 4.5 cm}
    \begin{paracol}{2}
}{
    \switchcolumn \raggedleft \secondColumn
    \end{paracol}
    \endonecolentry
} % new environment for two column entries

\newenvironment{threecolentry}[3][]{
    \onecolentry
    \def\thirdColumn{#3}
    \setcolumnwidth{, \fill, 4.5 cm}
    \begin{paracol}{3}
    {\raggedright #2} \switchcolumn
}{
    \switchcolumn \raggedleft \thirdColumn
    \end{paracol}
    \endonecolentry
} % new environment for three column entries

\newenvironment{header}{
    \setlength{\topsep}{0pt}\par\kern\topsep\centering\linespread{1.5}
}{
    \par\kern\topsep
} % new environment for the header

\newcommand{\placelastupdatedtext}{% \placetextbox{<horizontal pos>}{<vertical pos>}{<stuff>}
  \AddToShipoutPictureFG*{% Add <stuff> to current page foreground
    \put(
        \LenToUnit{\paperwidth-2 cm-0 cm+0.05cm},
        \LenToUnit{\paperheight-1.0 cm}
    ){\vtop{{\null}\makebox[0pt][c]{
        \small\color{gray}\textit{Last updated in August 2026}\hspace{\widthof{Last updated in August 2026}}
    }}}%
  }%
}%

% save the original href command in a new command:
\let\hrefWithoutArrow\href




\begin{document}
    \placelastupdatedtext

    \newcommand{\AND}{\unskip
        \cleaders\copy\ANDbox\hskip\wd\ANDbox
        \ignorespaces
    }
    \newsavebox\ANDbox
    \sbox\ANDbox{$|$}

    \begin{header}
        \fontsize{22 pt}{22 pt}\selectfont Dimitry Grebenyuk

        \vspace{4 pt}

        \normalsize
        \mbox{Shenzhen, China}%
        \kern 5.0 pt%
        \AND%
        \kern 5.0 pt%
        \mbox{\hrefWithoutArrow{mailto:dimitrygrebenyuk@icloud.com}{dimitrygrebenyuk@icloud.com}}%
        \kern 5.0 pt%
        \AND%
        \kern 5.0 pt%
        \mbox{\hrefWithoutArrow{tel:+86-135-3004-2782}{+86-13530042782}}%
        \kern 5.0 pt%
        \AND%
        \kern 5.0 pt%
        \mbox{\hrefWithoutArrow{https://github.com/grebenyyk}{github.com/grebenyyk}}%
        \kern 5.0 pt%
        \AND%
        \kern 5.0 pt%
        \mbox{\hrefWithoutArrow{https://orcid.org/0000-0002-4723-8564}{ORCID}}
    \end{header}

    \vspace{5 pt - 0.3 cm}

    \section{Research Interests}
    \vspace{0.04 cm}
    \begin{onecolentry}
    Lanthanide coordination chemist and materials scientist (PhD, Lomonosov Moscow State University), currently Senior Lecturer at Shenzhen MSU-BIT University. My research focuses on developing rare-earth metal–organic frameworks and polynuclear clusters, combining conventional crystallography with total-scattering analysis and machine learning–driven structural modeling (DFT, molecular simulations) to resolve their structural complexity. \newline I lead inorganic chemistry teaching and supervise research on functional materials.
    \end{onecolentry}

    \section{Education}
        \begin{twocolentry}{
            2018 – 2022
        }
            \textbf{Lomonosov Moscow State University}, PhD, Inorganic Chemistry\newline
            Thesis: «Coordination polymers and polynuclear clusters based on aliphatic rare earth carboxylates»
            \end{twocolentry}

        \vspace{0.03 cm}

        \begin{twocolentry}{
            2012 – 2018
        }
            \textbf{Lomonosov Moscow State University}, MSc \textit{cum laude}, Chemistry \end{twocolentry}
        \vspace{0.04 cm}

    \section{Academic Appointments}
        \begin{twocolentry}{
            2023 – Present
        }
            \textbf{Senior Lecturer}, Faculty of Materials Science, Shenzhen MSU-BIT University\newline
            {\small Instructor and research supervisor in chemistry and materials science (previously Teaching Assistant, 2019–2023).}
            \end{twocolentry}

        \vspace{0.03 cm}

        \begin{twocolentry}{
            2018 – 2023
        }
            \textbf{Doctoral Researcher \& Research Engineer}, Department of Chemistry, Lomonosov Moscow State University\newline
            {\small Doctoral research on lanthanide coordination polymers; concurrent diffraction-support engineering role.}
            \end{twocolentry}
        \vspace{0.04 cm}

    \section{Selected Publications}
    \vspace{0.04 cm}
    \begin{onecolentry}
        \begin{itemize}[
            topsep=0.03 cm,
            parsep=0.05 cm,
            partopsep=0pt,
            itemsep=0pt,
            leftmargin=0 cm + 10pt
        ]
            \item \href{https://www.sciencedirect.com/science/article/pii/S0022459626002811}{\textbf{Grebenyuk, D.}, Sabitova, I., Long, H., Gavrilkin, S., Tsvetkov, A., Lyssenko, K., Tsymbarenko, D.
            Solid-state transformation and structural diversity in lanthanide 1D coordination polymers with bulky cyclohexanecarboxylate ligands.
            \textit{J. Solid State Chem.} \textbf{2026}, \textit{360}, 126085.}

            \item \href{https://pubs.acs.org/doi/10.1021/acsomega.3c07906}{\textbf{Grebenyuk, D.}, Shaulskaya, M., Shevchenko, A., Zobel, M., Tedeeva, M., Kustov, A., Sadykov, I., Tsymbarenko, D.
            Tuning the Cerium-Based Metal–Organic Framework Formation by Template Effect and Precursor Selection.
            \textit{ACS Omega} \textbf{2023}, \textit{8} (50), 48394–48404.}

            \item \href{https://www.mdpi.com/2073-4360/14/16/3328}{\textbf{Grebenyuk, D.}, Zobel, M., Tsymbarenko, D.
            Partially Ordered Lanthanide Carboxylates with a Highly Adaptable 1D Polymeric Structure.
            \textit{Polymers} \textbf{2022}, \textit{14} (16), 3328.}

            \item \href{https://doi.org/10.1107/S1600576722005878}{Tsymbarenko, D., \textbf{Grebenyuk, D.}, Burlakova, M., Zobel, M.
            Quick and robust PDF data acquisition using a laboratory single-crystal X-ray diffractometer for study of polynuclear lanthanide complexes in solid form and in solution.
            \textit{J. Appl. Cryst.} \textbf{2022}, \textit{55}, 890–900.}

            \item \href{https://pubs.acs.org/doi/10.1021/acs.inorgchem.1c00581}{\textbf{Grebenyuk, D.}, Zobel, M., Polentarutti, M., Ungur, L., Kendin, M., Zakharov, K., Degtyarenko, P., Vasiliev, A., Tsymbarenko, D.
            A Family of Lanthanide Hydroxo Carboxylates with 1D Polymeric Topology and Ln\textsubscript{4} Butterfly Core Exhibits Switchable Supramolecular Arrangement.
            \textit{Inorg. Chem.} \textbf{2021}, \textit{60} (11), 8049–8061.}
        \end{itemize}
        \vspace{0.06cm}
        \textit{Full publication list (h-index 7, $\sim$200 citations):} \href{https://scholar.google.com/citations?user=pNWZLJoAAAAJ&hl=en}{scholar.google.com}
    \end{onecolentry}

    \section{Teaching \& Supervision}
    \vspace{0.04 cm}
    \begin{onecolentry}
    \textbf{Courses developed and taught}
    \begin{highlights}
        \item Designed and delivered the undergraduate course \textit{Introduction to Chemistry}; authored the accompanying textbook, published by MSU-BIT.
        \item Teach full undergraduate courses: \textit{General Chemistry and Chemistry of Elements} (seminars and laboratory sessions) and \textit{Modern Inorganic Chemistry} (seminars).
        \item Deliver an annual lecture on the \textit{History of Materials Science} for master's students.
    \end{highlights}

    \vspace{0.06cm}
    \textbf{Student supervision}
    \begin{highlights}
        \item Supervised 23 bachelor's and master's thesis projects on metal–organic frameworks and luminescent materials; 16 students have continued to graduate programs.
        \item Alina Zhumabekova's thesis «Sorption of molecular iodine in lanthanide terephthalate metal–organic frameworks» was commended by the State Degree Committee overseeing the defense.
    \end{highlights}

    \vspace{0.06cm}
    \textbf{Recognition:} Second Prize, MSU-BIT University Competition for Pedagogical Innovations (2025), as a team member.
    \end{onecolentry}

    \section{Invited Talks \& Conferences}
    \vspace{0.04 cm}
    \begin{onecolentry}
        \begin{itemize}
            \item \textbf{Beijing Conference on Molecular Sciences}, 2025, Beijing — \textit{``Molecular clusters of f-elements: from synthesis to deep learning insights''} (invited)
            \item \textbf{Peking University International Forum on ``AI + New Materials''}, 2024, Shenzhen — \textit{``Deep learning insight into molecular clusters nuclearity via pair distribution function''} (invited)
        \end{itemize}
    \end{onecolentry}

    \section{Professional Service}
    \vspace{0.04 cm}
    \begin{onecolentry}
    Peer reviewer for \textit{Frontiers in Nuclear Engineering}.
    \end{onecolentry}

    \section{Software \& Data}
    \vspace{0.04 cm}
    \begin{onecolentry}
        \textbf{PDF-NN} — neural network for predicting the nuclearity of lanthanide clusters from pair distribution function (PDF) data. Source code, trained models, and accompanying dataset publicly released under the MIT license.\newline
        Code \& models: \href{https://github.com/grebenyyk/PDF-NN}{github.com/grebenyyk/PDF-NN} \quad Data: \href{https://doi.org/10.5281/zenodo.19173873}{Zenodo, DOI 10.5281/zenodo.19173873}

        \vspace{0.06cm}
        \textbf{OEIS A000602} — contributed the higher-order correction terms in the asymptotic formula for the number of alkane structural isomers, accepted into the OEIS entry.\newline
        Entry: \href{https://oeis.org/A000602}{oeis.org/A000602} \quad Code \& write-up: \href{https://github.com/grebenyyk/alkane-isomers}{github.com/grebenyyk/alkane-isomers}

        \vspace{0.06cm}
        \textbf{OEIS A397444} — derived the closed-form expression for the anion excess of saturated rock-salt (NaCl) nanocubes, $a(n) = 3n^2 + 1 - (-1)^n$, as an exact surface-stoichiometry descriptor for ionic nanocrystals.\newline
        Entry: \href{https://oeis.org/A397444}{oeis.org/A397444} \quad Code \& write-up: \href{https://github.com/grebenyyk/np-stoichiometry}{github.com/grebenyyk/np-stoichiometry}
    \end{onecolentry}

    % Uncomment and fill in when an independent grant is secured, or to list
    % participant roles on funded projects (funder, title, period, role).
    % \section{Grants \& Funding}
    % \vspace{0.04 cm}
    % \begin{onecolentry}
    %     \begin{itemize}
    %         \item ...
    %     \end{itemize}
    % \end{onecolentry}

    \section{Skills \& Expertise}
    \vspace{0.04 cm}
    \begin{onecolentry}
    \textbf{Research:} Lanthanide coordination chemistry, Metal–organic frameworks,
    Crystallography (SC-XRD, PXRD), Total scattering and PDF analysis,
    Synchrotron beamtime (Diamond, Elettra, Kurchatov)

    \vspace{0.06cm}
    \textbf{Computing:} Python, Machine learning for structure prediction from PDF data, DFT \& molecular simulations (VASP, Quantum Espresso, RASPA), High-performance computing, Crystallographic software (SHELX, Olex2)

    \vspace{0.06cm}
    \textbf{Languages:} Russian (native), English (fluent), Mandarin Chinese (HSK4)
    \end{onecolentry}
\end{document}
